% ====================================================================
% signalChain.tex
% --------------------------------------------------------------------
% Block diagram of the sensor channel, used as the key figure on the
% title page. Drawn with plain TikZ.
% ====================================================================

\begin{figure}[H]
\centering
\begin{tikzpicture}[
  block/.style={draw, thick, minimum width=2.6cm, minimum height=1.3cm,
                align=center, font=\small},
  wanted/.style={block, fill=linkblue!12},
  >=latex
]
  \node[block]  (sensor) at (0,0)    {Sensor\\amplifier};
  \node[wanted] (filter) at (3.8,0)  {Anti-aliasing\\filter};
  \node[block]  (adc)    at (7.6,0)  {12-bit\\ADC};
  \node[block]  (mcu)    at (11.4,0) {Firmware};

  \draw[->, thick] (sensor) -- (filter);
  \draw[->, thick] (filter) -- (adc);
  \draw[->, thick] (adc)    -- (mcu);

  \node[font=\footnotesize, below=2pt] at (sensor.south) {DC to \qty{500}{\hertz}};
  \node[font=\footnotesize, below=2pt] at (filter.south) {$f_0 = \qty{1}{\kilo\hertz}$};
  \node[font=\footnotesize, below=2pt] at (adc.south)    {$f_s = \qty{10}{\kilo\hertz}$};

  \draw[decorate, decoration={brace, amplitude=6pt}, thick]
    ([yshift=8pt]filter.north west) -- ([yshift=8pt]filter.north east)
    node[midway, above=8pt, font=\footnotesize] {covered by this note};
\end{tikzpicture}
\caption{The sensor channel, with the stage this note covers shaded.}
\label{fig:signalChain}
\end{figure}
